\documentclass[11pt]{article}
\usepackage[T1]{fontenc}
\usepackage[utf8]{inputenc}
\usepackage{lmodern}
\usepackage{amsmath,amssymb,amsthm,mathtools}
\usepackage{booktabs,array}
\usepackage[margin=1in]{geometry}
\usepackage{microtype}
\usepackage{xcolor}
\usepackage{hyperref}
\usepackage{xurl}
\hypersetup{
  colorlinks=true,
  linkcolor=blue!55!black,
  citecolor=blue!55!black,
  urlcolor=blue!55!black,
  pdftitle={Paley two-edge switching for proper transposed sesqui arrays},
  pdfauthor={Carptopus}
}
\newtheorem{theorem}{Theorem}[section]
\newtheorem{lemma}[theorem]{Lemma}
\newtheorem{proposition}[theorem]{Proposition}
\newtheorem{corollary}[theorem]{Corollary}
\theoremstyle{definition}
\newtheorem{definition}[theorem]{Definition}
\theoremstyle{remark}
\newtheorem{remark}[theorem]{Remark}
\DeclareMathOperator{\GF}{GF}
\newcommand{\F}{\mathbb F}
\newcommand{\SA}{SA^{\mathsf T}}
\newcommand{\calA}{\mathcal A}
\newcommand{\one}{\mathbf 1}

\title{Paley two-edge switching for proper transposed sesqui arrays}
\author{Carptopus\\\href{mailto:carptopus@163.com}{\texttt{carptopus@163.com}}}
\date{22 August 2026\\\small Version v0.3-beta}

\begin{document}
\maketitle

\begin{abstract}
Let \(q\equiv3\pmod4\) be an odd prime power.  We construct, for every
\(q\geq11\), a proper transposed sesqui array
\[
 \SA\!\left(
 2q,\frac{q+1}{2},-,\frac{q-1}{2},\frac{q+1}{2}:q\times(q+1)
 \right).
\]
The construction starts from a Paley matching between the nonzero squares
and nonsquares of \(\F_q\), then crosses two matching edges.  Five quadratic
character conditions preserve the column--column and row--column
concurrences while a nonzero second-moment defect forces the row--row
concurrence to be nonconstant.  An exact character count reduces the
existence of suitable switching parameters to the Hasse bound for a family
of elliptic curves.  The bound covers \(q\geq19\), and the remaining case
\(q=11\) has an explicit witness.
\end{abstract}

\paragraph{Status.}
This is a proof-complete initial draft.  Its symbolic identities, finite-field
instances, and boundary cases have passed implementation-separated internal
checks.  External mathematical review and a definitive priority assessment
remain open, so the novelty statement below is qualified.

\section{Introduction}
Row--column designs are binary arrays whose rows and columns are regarded as
blocks on a common symbol set.  Triple arrays require constant row--row,
column--column, and row--column concurrence.  A transposed sesqui array keeps
the latter two conditions but does not require constant row--row concurrence;
it is \emph{proper} when row--row concurrence is in fact nonconstant
\cite{jager-enumeration}.

For odd prime powers \(q\equiv3\pmod4\) with \(q>3\), the classical Paley construction
gives triple arrays with parameters
\[
 TA\!\left(
 2q,\frac{q+1}{2},\frac{q+1}{2},\frac{q-1}{2},
 \frac{q+1}{2}:q\times(q+1)
 \right)
\]
\cite{preece-paley,nilson-intercalates}.  Thus the dimensions, replication,
and concurrence parameters studied here are not new.  Our objective is to
retain the Paley column structure while replacing the constant row
concurrence by a provably nonconstant one.

The construction is local.  A standard Paley matching pairs every nonzero
square with a nonsquare.  We cross two of its edges, subject to five quadratic
character conditions.  This changes only four endpoints, but it has three
global consequences: the sum map remains two-to-one away from zero, the
column design remains Paley, and the row block acquires a nonzero moment
defect.  The first two properties give a transposed sesqui array; the third
proves that it is proper.

The nonlocal part of the argument is existence.  For a fixed matching
parameter \(\alpha\), the number of valid second edges is expressed exactly as
a cubic quadratic-character sum.  The associated curve is elliptic, and the
Hasse bound makes the count positive for every \(q\geq19\).  A direct witness
closes \(q=11\).  No finite computation is used to infer the infinite family.

The smallest instance has \((v,r,c)=(22,11,12)\).  The 2025 enumeration and
construction table of J\"ager et al. marks its proper transposed-sesqui entry
as \texttt{?,T}: a same-parameter triple array was known, while their search
and listed constructions did not produce a proper example
\cite[Table~2]{jager-enumeration}.  To the best of our knowledge, the Paley
two-edge switching construction and the resulting proper family have not
appeared in the row--column design literature.

\section{Definitions and Paley preliminaries}
Let \(A\) be an \(r\times c\) array on a symbol set \(\Omega\) of size \(v\).
Write \(R_i\) and \(C_j\) for the symbol supports of row \(i\) and column
\(j\).

\begin{definition}
The array \(A\) is \emph{binary} if no symbol occurs twice in one row or one
column, and \emph{equireplicate} with replication \(e\) if every symbol occurs
in exactly \(e\) cells.  It is a transposed sesqui array if there are constants
\(\lambda_{CC}\) and \(\lambda_{RC}\) such that
\[
 |C_j\cap C_{j'}|=\lambda_{CC}\quad(j\ne j'),
 \qquad
 |R_i\cap C_j|=\lambda_{RC}\quad\text{for all }i,j.
\]
We write
\[
 \SA(v,e,-,\lambda_{CC},\lambda_{RC}:r\times c).
\]
It is \emph{proper} if \(|R_i\cap R_{i'}|\) is not constant over distinct
rows.
\end{definition}

Let \(F=\F_q\), where throughout \(q\equiv3\pmod4\) is an odd prime power.
Let \(\chi:F\to\{0,\pm1\}\) be the quadratic character, with \(\chi(0)=0\),
and put
\[
 Q=\{x^2:x\in F^*\},\qquad N=F^*\setminus Q,
 \qquad N_0=N\cup\{0\}.
\]
Then \(\chi(-1)=-1\), and \(Q\sqcup N_0=F\).

\begin{lemma}[Paley intersections]
\label{lem:paley}
For every \(h\in F^*\),
\[
 |Q\cap(Q+h)|=\frac{q-3}{4},\qquad
 |N_0\cap(N_0+h)|=\frac{q+1}{4}.
\]
In particular, \(Q\) is a
\((q,(q-1)/2,(q-3)/4)\) difference set in the additive group of \(F\),
and \(N_0\) is its complement.
\end{lemma}

\begin{proof}
For \(h\ne0\), the standard quadratic sum
\[
 \sum_{x\in F}\chi\bigl(x(x-h)\bigr)=-1
\]
follows by completing the square.  Indeed, after scaling by \(h\), it reduces
to \(\sum_t\chi(t^2-1)=-1\); the equation \(y^2=t^2-1\) has \(q-1\)
solutions because \((t-y)(t+y)=1\), and counting its solutions by \(t\)
gives the stated sum.  Hence
\[
 \frac14\sum_{x\in F}(1+\chi(x))(1+\chi(x-h))=\frac{q-1}{4}.
\]
This product is the indicator of \(Q\cap(Q+h)\), except at \(x=0,h\).
The two false contributions total
\[
 \frac{1+\chi(-h)}4+\frac{1+\chi(h)}4=\frac12,
\]
because \(\chi(-1)=-1\).  Thus \(|Q\cap(Q+h)|=(q-3)/4\).
Since \(N_0=F\setminus Q\), taking complements gives
\(|N_0\cap(N_0+h)|=(q+1)/4\).
\end{proof}

\section{A switched Paley matching}
We first isolate the conditions on the two crossed edges.

\begin{proposition}[Two-edge switch]
\label{prop:switch}
Suppose \(\alpha,r\in F\) satisfy
\begin{equation}
 \chi(\alpha)=\chi(1+\alpha)=-1,
 \label{eq:alpha-conditions}
\end{equation}
and
\begin{equation}
 \chi(r)=\chi(1+\alpha r)=\chi(r+\alpha)=1.
 \label{eq:r-conditions}
\end{equation}
Define \(\pi_0:F\to F\) by
\begin{equation}
 \pi_0(0)=0,\qquad
 \pi_0(u)=\alpha u\ (u\in Q),\qquad
 \pi_0(v)=\alpha^{-1}v\ (v\in N).
 \label{eq:baseline-involution}
\end{equation}
Replace the two pairs
\[
 \{1,\alpha\},\qquad \{r,\alpha r\}
\]
by
\[
 \{1,\alpha r\},\qquad \{r,\alpha\},
\]
and leave every other pair of \(\pi_0\) unchanged.  The resulting map
\(\pi\) is an involution.  Moreover, the map
\begin{equation}
 s:F\longrightarrow F,\qquad s(b)=b+\pi(b),
 \label{eq:sum-map}
\end{equation}
has one singleton fibre \(\{0\}\) over zero and \((q-1)/2\) two-element
fibres over distinct nonzero values.  Each two-element fibre contains one
square and one nonsquare.
\end{proposition}

\begin{proof}
Equation \eqref{eq:baseline-involution} pairs each member of \(Q\) with a
member of \(N\), and conversely, because \(\chi(\alpha)=-1\).  The condition
\(r=1\) is impossible: it would make
\(\chi(1+\alpha r)=\chi(1+\alpha)=-1\), contrary to
\eqref{eq:r-conditions}.  Hence \(1,r,\alpha,\alpha r\) are four distinct
nonzero endpoints, and the crossed pairs still join \(Q\) to \(N\).

Every unchanged pair has sum \((1+\alpha)u\) for some
\(u\in Q\setminus\{1,r\}\).  These sums are distinct and lie in \(N\) by
\eqref{eq:alpha-conditions}.  The two new sums are
\[
 1+\alpha r,\qquad r+\alpha,
\]
which lie in \(Q\) by \eqref{eq:r-conditions}; hence they cannot collide with
an unchanged sum.  They do not collide with one another because
\[
 (1+\alpha r)-(r+\alpha)=(1-\alpha)(1-r)\ne0.
\]
Together with the fixed point \(0\), this proves all assertions about the
fibres of \(s\).
\end{proof}

Let
\begin{equation}
 D=s(F)=\{b+\pi(b):b\in F\}.
 \label{eq:D-definition}
\end{equation}
Proposition~\ref{prop:switch} gives \(|D|=(q+1)/2\).  Relative to the
unswitched sum set \(N_0\), it has the explicit form
\begin{equation}
 D=\bigl(N_0\setminus\{1+\alpha,r(1+\alpha)\}\bigr)
 \cup\{1+\alpha r,r+\alpha\}.
 \label{eq:D-switch}
\end{equation}

\section{The cellwise array construction}
The next step records the ordering, including the zero fibre and the extra
column.

\begin{theorem}[Array from the switched matching]
\label{thm:cellwise}
Under the hypotheses of Proposition~\ref{prop:switch}, there is a binary
equireplicate transposed sesqui array
\[
 \SA\!\left(
 2q,\frac{q+1}{2},-,\frac{q-1}{2},\frac{q+1}{2}:q\times(q+1)
 \right).
\]
Its row--row concurrence is nonconstant if and only if the additive
intersection numbers \(|D\cap(D+h)|\), \(h\in F^*\), are nonconstant.
\end{theorem}

\begin{proof}
Take the symbol set \(\Omega=F\times\{0,1\}\), index rows by \(a\in F\), and
index columns by \(F\cup\{\infty\}\).  Define
\[
 \ell(b)=
 \begin{cases}
 0,&\pi(b)\in Q,\\
 1,&\pi(b)\in N_0,
 \end{cases}
\]
and, for a finite column \(y\), put \(b=y-a\) and
\begin{equation}
 A(a,y)=\bigl(a+b+\pi(b),\ell(b)\bigr).
 \label{eq:array-cell}
\end{equation}
In the additional column set
\begin{equation}
 A(a,\infty)=(a,0).
 \label{eq:infinity-column}
\end{equation}

Fix a row \(a\).  Each nonzero fibre of \(s\) contains a square and a
nonsquare, which receive different labels because \(\pi\) interchanges the
two character classes.  Thus every \(d\in D\setminus\{0\}\) contributes both
\((a+d,0)\) and \((a+d,1)\).  The zero fibre is \(b=0\): its finite cell is
\((a,1)\), while \eqref{eq:infinity-column} supplies \((a,0)\).  Therefore
the row support is exactly
\begin{equation}
 R_a=(a+D)\times\{0,1\}.
 \label{eq:row-support}
\end{equation}
In particular, every row is binary.

Fix a finite column \(y\).  As \(a\) runs through \(F\), the element
\(b=y-a\) runs through \(F\), and \(\pi(b)\) also runs through \(F\).  Since
the first coordinate in \eqref{eq:array-cell} is \(y+\pi(b)\), the column
support is
\begin{equation}
 C_y=((y+Q)\times\{0\})\cup((y+N_0)\times\{1\}),
 \label{eq:finite-column-support}
\end{equation}
while
\begin{equation}
 C_\infty=F\times\{0\}.
 \label{eq:infinity-support}
\end{equation}
These columns are binary.  A symbol with label zero occurs in \(|Q|\) finite
columns and the infinity column; a symbol with label one occurs in \(|N_0|\)
finite columns.  Hence every symbol has replication \((q+1)/2\).

For two finite columns, Lemma~\ref{lem:paley} and the disjoint labels give
\[
 |C_y\cap C_{y'}|
 =\frac{q-3}{4}+\frac{q+1}{4}=\frac{q-1}{2}.
\]
A finite column meets \(C_\infty\) in its \(|Q|=(q-1)/2\) label-zero
symbols, so column concurrence is constant.  Finally, \(Q\sqcup N_0=F\)
implies
\[
 |R_a\cap C_y|
 =|(a+D)\cap(y+Q)|+|(a+D)\cap(y+N_0)|
 =|D|=\frac{q+1}{2},
\]
and the same value follows from \eqref{eq:infinity-support}.  This proves the
transposed-sesqui conditions.

For \(h\ne0\), equation \eqref{eq:row-support} gives
\begin{equation}
 |R_a\cap R_{a+h}|=2|D\cap(D+h)|.
 \label{eq:row-concurrence}
\end{equation}
The final assertion follows.
\end{proof}

\section{A moment obstruction to constant row concurrence}
We now show that the two-edge switch forces the array in
Theorem~\ref{thm:cellwise} to be proper.

\begin{lemma}[Difference-set moment identity]
\label{lem:moment}
Let \(E\subseteq F\), where \(q>3\) is odd, and set
\[
 k=|E|,\qquad S_1(E)=\sum_{x\in E}x,
 \qquad S_2(E)=\sum_{x\in E}x^2.
\]
If \(|E\cap(E+h)|\) is constant for \(h\in F^*\), then
\begin{equation}
 kS_2(E)-S_1(E)^2=0
 \label{eq:moment-necessary}
\end{equation}
in \(F\).
\end{lemma}

\begin{proof}
Constant nonzero translation intersections make \(E\) a difference set in
the additive group of \(F\); write its difference multiplicity as \(\lambda\).
Summing the square of every ordered difference gives
\[
 \sum_{x,y\in E}(x-y)^2
 =2kS_2(E)-2S_1(E)^2
 =\lambda\sum_{h\in F^*}h^2.
\]
The last field sum is zero.  For example, choose \(c\in F^*\) with
\(c^2\ne1\); multiplication by \(c\) permutes \(F^*\), so the sum equals
\(c^2\) times itself.  Since the characteristic is odd, division by two
proves \eqref{eq:moment-necessary}.
\end{proof}

\begin{proposition}[Nonzero switched moment]
\label{prop:proper}
The block \(D\) in \eqref{eq:D-switch} satisfies
\begin{equation}
 |D|S_2(D)-S_1(D)^2=-\alpha(r-1)^2\ne0.
 \label{eq:moment-defect}
\end{equation}
Consequently, the array in Theorem~\ref{thm:cellwise} is proper.
\end{proposition}

\begin{proof}
The removed and inserted elements in \eqref{eq:D-switch} have the same sum:
\[
 (1+\alpha)+r(1+\alpha)=(1+\alpha r)+(r+\alpha).
\]
Thus the first moment does not change.  Direct expansion gives the change in
the second moment:
\begin{align*}
 \Delta S_2
 &=(1+\alpha r)^2+(r+\alpha)^2
   -(1+\alpha)^2-r^2(1+\alpha)^2\\
 &=-2\alpha(r-1)^2.
\end{align*}
The original block \(N_0\) is a difference set, so its moment defect is zero
by Lemma~\ref{lem:moment}.  In the field \(F\), the integer
\(|D|=(q+1)/2\) equals \(1/2\), because \(q=0\) in \(F\).  Hence the new
defect is
\[
 |D|\Delta S_2=-\alpha(r-1)^2,
\]
which is nonzero.  Lemma~\ref{lem:moment} shows that the translation
intersections of \(D\) cannot be constant; now use
\eqref{eq:row-concurrence}.
\end{proof}

\section{Existence of switching parameters}
It remains to prove that \(\alpha,r\) satisfying
\eqref{eq:alpha-conditions}--\eqref{eq:r-conditions} always exist in the
claimed range.

\begin{lemma}[The first parameter]
\label{lem:alpha-count}
Let
\[
 \calA=\{\alpha\in F:\chi(\alpha)=\chi(1+\alpha)=-1\}.
\]
Then
\begin{equation}
 |\calA|=\frac{q-3}{4}.
 \label{eq:alpha-count}
\end{equation}
\end{lemma}

\begin{proof}
The standard quadratic-character sum is
\[
 \sum_{\alpha\in F}\chi\bigl(\alpha(1+\alpha)\bigr)=-1.
\]
Expanding
\((1-\chi(\alpha))(1-\chi(1+\alpha))/4\) and summing therefore gives
\((q-1)/4\).  This expression is the indicator of membership in \(\calA\)
except at the zeros.  At \(\alpha=0\) it contributes zero, while at
\(\alpha=-1\) it contributes \(1/2\), although \(-1\notin\calA\).  Removing
this false contribution proves \eqref{eq:alpha-count}.
\end{proof}

Fix \(\alpha\in\calA\), and define
\begin{align}
 C(\alpha)
 &=\#\{r\in F:\chi(r)=\chi(1+\alpha r)=\chi(r+\alpha)=1\},
 \label{eq:C-definition}\\
 T(\alpha)
 &=\sum_{r\in F}\chi\bigl(r(1+\alpha r)(r+\alpha)\bigr).
 \label{eq:T-definition}
\end{align}

\begin{lemma}[Exact character count]
\label{lem:character-count}
For every \(\alpha\in\calA\),
\begin{equation}
 C(\alpha)=
 \frac{q-3+T(\alpha)-4\chi(\alpha-1)}{8}.
 \label{eq:exact-character-count}
\end{equation}
\end{lemma}

\begin{proof}
Temporarily sum the product of the three positive-character indicators:
\[
 \frac18\sum_{r\in F}
 (1+\chi(r))(1+\chi(1+\alpha r))(1+\chi(r+\alpha)).
\]
 The three linear character sums vanish.  For a quadratic polynomial with
 nonzero leading coefficient and nonzero discriminant, completing the square
 gives
 \(\sum_x\chi(Ax^2+Bx+C)=-\chi(A)\).  Since
 \(\alpha\notin\{0,\pm1\}\), this identity gives the three pairwise sums
\begin{align*}
 \sum_r\chi(r(1+\alpha r))&=1,\\
 \sum_r\chi(r(r+\alpha))&=-1,\\
 \sum_r\chi((1+\alpha r)(r+\alpha))&=1.
\end{align*}
Together with \eqref{eq:T-definition}, the uncorrected total is
\((q+1+T(\alpha))/8\).

The zero \(r=0\) contributes nothing to that expanded expression.  The two
other zeros, \(r=-1/\alpha\) and \(r=-\alpha\), are distinct and each
contributes falsely
\[
 \frac{1+\chi(\alpha-1)}{4}.
\]
Indeed, the remaining nonzero character is \(\chi(\alpha-1)\) at each
point, while the other nonzero factor has character one.
Subtracting both contributions gives
\eqref{eq:exact-character-count}.
\end{proof}

\begin{lemma}[Elliptic-curve bound]
\label{lem:hasse}
For every \(\alpha\in\calA\),
\begin{equation}
 |T(\alpha)|\leq2\sqrt q.
 \label{eq:hasse-bound}
\end{equation}
Consequently,
\begin{equation}
 C(\alpha)\geq\frac{q-7-2\sqrt q}{8}.
 \label{eq:C-lower}
\end{equation}
\end{lemma}

\begin{proof}
Consider
\begin{equation}
 E_\alpha:\quad z^2=r(1+\alpha r)(r+\alpha).
 \label{eq:elliptic-curve}
\end{equation}
The roots \(0,-1/\alpha,-\alpha\) are pairwise distinct because
\(\alpha\notin\{0,\pm1\}\).  Hence the affine cubic is nonsingular.  Its
projective closure
\[
 Z^2W=R(W+\alpha R)(R+\alpha W)
\]
has the unique point \([0:1:0]\) at infinity, and the partial derivative with
respect to \(W\) is nonzero there.  Thus \(E_\alpha\) is an elliptic curve
over \(F\).

For each \(r\in F\), the number of affine \(z\)-values is
\(1+\chi(r(1+\alpha r)(r+\alpha))\).  Including the point at infinity gives
\[
 \#E_\alpha(F)=q+1+T(\alpha).
\]
The Hasse bound for elliptic curves over finite fields
\cite[Chapter~V]{silverman} proves \eqref{eq:hasse-bound}.  Substituting
\(T(\alpha)\geq-2\sqrt q\) and \(\chi(\alpha-1)\leq1\) into
\eqref{eq:exact-character-count} proves \eqref{eq:C-lower}.
\end{proof}

\begin{proposition}[Parameter existence]
\label{prop:parameter-existence}
If \(q\equiv3\pmod4\) is an odd prime power and \(q\geq19\), then there are
\(\alpha,r\in F\) satisfying
\eqref{eq:alpha-conditions}--\eqref{eq:r-conditions}.
\end{proposition}

\begin{proof}
Lemma~\ref{lem:alpha-count} gives \(\calA\ne\varnothing\).  For any
\(\alpha\in\calA\), the right-hand side of \eqref{eq:C-lower} is positive
when \(q\geq19\).  Since \(C(\alpha)\) is an integer, a valid \(r\) exists.
As observed in Proposition~\ref{prop:switch}, its conditions automatically
exclude \(r=1\).
\end{proof}

\section{The boundary instance and the main theorem}
For \(q=11\), take \(\alpha=7\) and \(r=5\).  The involution consists of the
fixed point \(0\) and the pairs
\[
 \{1,2\},\quad\{5,7\},\quad\{3,10\},\quad\{4,6\},\quad\{9,8\}.
\]
All five character conditions hold.  Its sum set is
\begin{equation}
 D=\{0,1,2,3,6,10\}\subset\F_{11},
 \label{eq:q11-D}
\end{equation}
and \eqref{eq:moment-defect} equals
\[
 -7(5-1)^2=9\ne0\pmod{11}.
\]
Equations \eqref{eq:array-cell} and \eqref{eq:infinity-column} therefore give
a proper
\[
 \SA(22,6,-,5,6:11\times12).
\]

\begin{theorem}[Main theorem]
\label{thm:main}
For every odd prime power \(q\equiv3\pmod4\) with \(q\geq11\), there exists
a proper transposed sesqui array
\[
 \SA\!\left(
 2q,\frac{q+1}{2},-,\frac{q-1}{2},\frac{q+1}{2}:q\times(q+1)
 \right).
\]
\end{theorem}

\begin{proof}
The explicit parameters above cover \(q=11\).  Every remaining admissible
prime power is at least \(19\), so Proposition~\ref{prop:parameter-existence}
provides switching parameters.  Proposition~\ref{prop:switch},
Theorem~\ref{thm:cellwise}, and Proposition~\ref{prop:proper} then construct
the required proper array.
\end{proof}

\begin{remark}[Lower parameters]
The switching conditions themselves have no solution for \(q=3\) or \(7\):
\(\calA\) is empty for \(q=3\); for \(q=7\), the sole candidate is
\(\alpha=5\), and no \(r\) satisfies \eqref{eq:r-conditions}.  This is a
boundary of the present construction, not a nonexistence claim for all
proper arrays with those parameters.
\end{remark}

\section{Prior-work boundary}
The target parameter family overlaps substantially with earlier work.  The
Paley triple arrays of Preece, Wallis, and Yucas already have the same
dimensions, replication, column concurrence, and row--column concurrence
\cite{preece-paley}.  They have constant row concurrence; later work has
studied their Youden-square representations and intercalates
\cite{nilson-youden,nilson-intercalates}.  Difference-set constructions give
another infinite source of triple arrays, again with constant row concurrence
by definition \cite{nilson-cameron}.

The deletion--exchange construction of J\"ager et al. can transform a Youden
rectangle into a row--column design and automatically preserves a binary,
equireplicate column-balanced layer \cite[Construction~5.2 and
Theorem~5.3]{jager-youden}.  It does not automatically provide the
row--column concurrence and properness proved here.  Recent resolvable triple
arrays also include an infinite Paley subfamily, but remain triple arrays
\cite{gordeev-ohman}.

Intercalates in Paley arrays are an established subject
\cite{nilson-intercalates}.  We therefore do not claim the Paley support, the
target parameters, the character-sum method, or the Hasse bound as new.  The
contribution claimed here is the
specific two-edge switch satisfying
\eqref{eq:alpha-conditions}--\eqref{eq:r-conditions}, the moment obstruction
that makes its array proper, and the existence proof for every odd prime
power \(q\equiv3\pmod4\), \(q\geq11\).

The closest dated parameter anchor is the \((22,11,12)\) entry in the 2025
table of J\"ager et al. \cite[Table~2]{jager-enumeration}.  That table is
evidence of an open construction gap at the time of publication, not by
itself proof of global priority.  Accordingly, our novelty statement remains
``to the best of our knowledge.''

The general Youden deletion interface does not imply this construction.
J\"ager et al. also obtain several proper arrays through that interface, so
neither the abstract deletion--exchange operation nor properness after a
Youden deletion is claimed as new.  The \(q=11\) array constructed here has a
compatible Youden completion, as certified in the accompanying verification
package.  What remains specific to this paper is the uniform Paley
parameterization, the two-edge switch and character conditions, and the
all-\(q\) existence proof.  Whether every larger member has a compatible
Youden representation is not used and remains open here.

The externally verifiable public archive date for the first version is no
later than 21 August 2026; manuscript dates and local Git timestamps are not
used as independent priority evidence.

\section{Reproducibility}
The proof is independent of computation.  The accompanying repository
contains six complementary verification paths:
\begin{itemize}
\item a symbolic check of the first- and second-moment identities;
\item an exact check of the character-count formula over sample prime fields;
\item full cellwise array checks over sample prime fields;
\item an implementation-independent check over \(\GF(27)\); and
\item a generic finite-field audit covering positive and negative boundary
      cases; and
\item an implementation-separated verification of the frozen (q=11)
      Youden completion and deletion roundtrip, including a destructive control.
\end{itemize}
The corresponding scripts are
\begin{quote}
\ttfamily\small
verify\_general\_paley\_switch\_symbolic.py\\
verify\_general\_paley\_switch\_character\_count.py\\
verify\_general\_paley\_switch\_prime\_fields.py\\
verify\_general\_paley\_switch\_gf27\_independent.py\\
audit\_general\_paley\_switch\_generic\_fields.py\\
audit\_q11\_youden\_certificate\_independent.py
\end{quote}
\medskip
They test the formulas and finite instances; they are not used to extrapolate
the infinite conclusion in Theorem~\ref{thm:main}.

\section*{Declaration of AI-assisted research and writing}
This research program used OpenAI Codex extensively for candidate selection,
conjecture formation, finite-field derivations, adversarial proof checking,
literature-search assistance, verification-code generation, and manuscript
drafting and revision.  The retained claims are presented through
human-readable arguments and reproducible artifacts, but AI output is not
treated as an authority and independent expert review remains pending.  Codex
is a tool, not an author.  The named author, Carptopus, accepts responsibility
for the manuscript's accuracy, originality, citations, and integrity.

\section*{License}
\textcopyright\ 2026 Carptopus.  This manuscript is licensed under the
\href{https://creativecommons.org/licenses/by/4.0/}{Creative Commons
Attribution 4.0 International License (CC BY 4.0)}.  Reusers must provide
appropriate attribution and indicate whether changes were made.  No
warranties are given.

\begingroup
\small
\begin{thebibliography}{99}
\bibitem{gordeev-ohman}
A.~Gordeev and L.-D. \"Ohman,
Resolvable triple arrays,
\emph{Electron. J. Combin.} 33(2) (2026), P2.35,
\url{https://doi.org/10.37236/14977}.

\bibitem{jager-enumeration}
G.~J\"ager, K.~Markstr\"om, D.~Shcherbak, and L.-D. \"Ohman,
Enumeration and construction of row-column designs,
\emph{J. Combin. Des.} 33 (2025), 357--372,
\url{https://doi.org/10.1002/jcd.21991}.

\bibitem{jager-youden}
G.~J\"ager, K.~Markstr\"om, D.~Shcherbak, and L.-D. \"Ohman,
Small Youden rectangles, near Youden rectangles, and their connections to
other row-column designs,
\emph{Discrete Math. Theor. Comput. Sci.} 25(1) (2023), Paper No.~9,
\url{https://doi.org/10.46298/dmtcs.6754}.

\bibitem{nilson-cameron}
T.~Nilson and P.~J. Cameron,
Triple arrays from difference sets,
\emph{J. Combin. Des.} 25 (2017), 494--506,
\url{https://doi.org/10.1002/jcd.21569}.

\bibitem{nilson-intercalates}
T.~Nilson,
Intercalates in double and triple arrays,
\emph{J. Combin. Des.} 30 (2022), 135--151,
\url{https://doi.org/10.1002/jcd.21816}.

\bibitem{nilson-youden}
T.~Nilson and L.-D. \"Ohman,
Triple arrays and Youden squares,
\emph{Des. Codes Cryptogr.} 75 (2015), 429--451,
\url{https://doi.org/10.1007/s10623-014-9926-8}.

\bibitem{preece-paley}
D.~A. Preece, W.~D. Wallis, and J.~L. Yucas,
Paley triple arrays,
\emph{Australas. J. Combin.} 33 (2005), 237--246,
\url{https://ajc.maths.uq.edu.au/pdf/33/ajc_v33_p237.pdf}.

\bibitem{silverman}
J.~H. Silverman,
\emph{The Arithmetic of Elliptic Curves}, 2nd ed., Graduate Texts in
Mathematics 106, Springer, 2009,
\url{https://doi.org/10.1007/978-0-387-09494-6}.
\end{thebibliography}
\endgroup

\end{document}
